\section{Discussion}
\label{sec:discussion}

\subsection{What the Pilot Shows}

The pilot shows that, on matched compositional-evaluation prompts, the
PCA-projected last-token activation matrices of two architecturally
different frontier transformers admit a substantially better
orthogonal alignment than they do under random row permutation. This
holds at three pre-registered layer fractions, in two
architecturally-different model pairs (Phi-4 / Qwen3-4B in Phase~0,
Phi-4 / Qwen3-Next-80B in Phase~1), and within each of three distinct
compositional-evaluation domains in Phase~1 (math, logic,
political-economy). The result is real. It is reproducible. The
infrastructure that produces it is released for reuse and extension.

\subsection{What the Pilot Does Not Show}

The pilot does not show that the aligned PCA subspaces are
operator-invariant in the sense of Eq.~\eqref{eq:leakage}. The paper
defines that statistic and does not measure it. PCA top-$k$ components
identify high-variance directions; high-variance is not invariance.
Until the leakage statistic is computed (Study A,
\S\ref{sec:requiredwork}), the operator-theoretic interpretation of
the alignment remains a hypothesis the pilot motivates rather than
confirms.

The pilot does not show that the alignment isolates compositional
structure from lexical, template, syntactic, or tokenizer-induced
structure. The row-permutation null discriminates ``the models encode
the prompts at all'' from ``the models encode them identically up to
an orthogonal transformation in their respective PCA subspaces.'' Two
competent language models given the same English prompts will produce
matched-row activations more similar than scrambled-row activations
\emph{regardless} of any operator-theoretic substrate, because both
models encode prompt identity, lexical content, and template
structure. Until lexical, template, and syntax-scrambled baselines
have been run (Study B), the pilot result is consistent with both the
operator-theoretic interpretation and the much weaker
shared-language-geometry interpretation.

The pilot itself did not show that the alignment is robust to
held-out prompts; the original procedure fit and evaluated Procrustes
on the same probe set. The first-pass Study C reported in
\S\ref{sec:results-study-c} addresses this concern directly:
train/test gap is essentially zero ($\leq 0.002$) for random splits
and small ($\leq 0.09$) for cross-domain splits.

The pilot itself did not compute the promised CKA cross-validation;
the first-pass Study F-CKA reported in \S\ref{sec:results-study-f-cka}
addresses this commitment: CKA is high ($0.84$--$0.96$ on the full
set) and far above the permutation null. The alignment that Procrustes
identifies is corroborated by an alignment-free metric. The
methodological commitment is now satisfied for Phase~1.

The pilot, even augmented by Studies C and F-CKA, does \emph{not} show
that the alignment exceeds simple lexical and template baselines. The
first-pass Study B reported in \S\ref{sec:results-study-b} shows the
opposite: the LLM-LLM Procrustes residual is essentially indistinguishable
from one of the LLMs aligned against a TF-IDF representation of the
same prompts, and is worse than a TF-IDF/template-only alignment. The
follow-up Study B' (\S\ref{sec:results-study-b-prime}) reports the
result of explicitly residualizing lexical/template structure out of
the LLM activations and re-running the alignment. We treat the joint
B/B' finding as the binding evidentiary constraint of the present
paper.

The pilot does not show that the alignment survives the shared
training-distribution confound. Both models in each pair are trained
on heavily overlapping post-2024 web corpora. The aligned subspace
could plausibly be a residue of shared corpus statistics. Until a
training-distribution-disjoint control is run (Study D), the
substrate-independence reading is one of two possible readings of the
result, the other being shared-corpus residue.

The pilot does not show that the aligned subspace is causally
responsible for compositional behavior. The alignment is correspondence,
not causation. The Procrustes transformation $R^*$ predicts which
direction in $V_B$ corresponds to a given direction in $V_A$, but
that prediction has not been tested by intervention (Study E).

\subsection{What Would Make This an Operator-Theoretic Confirmation}

A subsequent paper that completes Studies A through G of
\S\ref{sec:requiredwork} and reports them jointly with the pilot
result would, if the results are positive, support the strong
operator-theoretic claim. We do not estimate the probability that the
strong claim survives the controls; the present pilot is consistent
with both the strong claim and several weaker alternatives, and the
controls are designed precisely to discriminate.

A version of the controls that finds:

\begin{itemize}
    \item Low invariance leakage on the recovered subspaces (Study A
        positive).
    \item Alignment quality substantially exceeding lexical and
        template baselines (Study B positive).
    \item Robust train/test alignment with cross-domain transfer
        (Study C positive).
    \item Alignment preserved under substantial training-distribution
        disjointness (Study D positive).
    \item Cross-model intervention predictions confirmed (Study E
        positive).
\end{itemize}

would constitute the publication claim that this paper's title was
originally written to assert. The honest current status is that the
claim is hypothesized, the test infrastructure exists, the pilot
result is consistent with the hypothesis, and the discriminating
controls remain to be run.

\subsection{Phase~2: Training-Distribution Control (the Binding Outstanding Confound)}

After the joint B/B' lexical control reported in
\S\ref{sec:initialstudies}, the binding outstanding confound is
shared training-distribution overlap (Study D in
\S\ref{sec:requiredwork}). The pilot is run on two model pairs whose
training distributions overlap substantially; until that overlap is
broken in a controlled way, even an alignment that survives lexical
residualization is consistent with shared-corpus residue rather than
substrate-independent computation. The control requires a model whose
training distribution is substantially disjoint from the Phase~1
pair's. Candidate controls:

\begin{itemize}
    \item A model trained on a non-web-pretrained corpus (e.g., a
        scientific-corpus model from the Galactica
        \cite{taylor2022galactica} or PMC-LLaMA family).
    \item A multilingual model whose training distribution is dominated
        by a non-English language (e.g., a French-pretrained
        decoder-only model).
    \item A small-scale model trained from scratch on a curated
        curriculum with documented disjointness from the pilot pair's
        training data.
\end{itemize}

If alignment survives the training-distribution swap, the
substrate-independence reading is meaningfully supported. If alignment
quality tracks training-distribution overlap (high alignment for
similar training, low alignment for disjoint training), the pilot
signal is reattributed to corpus-shared statistics and the
operator-theoretic interpretation is correspondingly weakened.

\subsection{Phase~3: Causal Intervention}

A positive COSI alignment, even after Phase~1 and Phase~2, is
correspondence evidence. The full operator-theoretic claim---that the
aligned subspace is the locus where the compositional reasoning is
\emph{performed}, not merely \emph{represented}---requires causal
intervention. The proposed test:

\begin{enumerate}
    \item Identify a direction $v_A \in V_A$ that, when ablated, degrades
        a specific compositional-reasoning behavior in $M_A$.
    \item Use the Procrustes transformation $R^*$ to map $v_A$ to its
        correspondent $v_B = R^{*-1} v_A \in V_B$.
    \item Predict the ablation effect of $v_B$ in $M_B$. Run the
        intervention.
    \item Compare predicted to observed effect.
\end{enumerate}

A successful Phase~3 result---ablation effects in $M_B$ that match those
predicted from interventions in $M_A$ via the Procrustes map---would
constitute the strongest possible empirical evidence for the
operator-theoretic framing of Proposition~\ref{prop:cosi}. We treat
Phase~3 as outside the scope of the present paper and as the natural
follow-up program.

\subsection{Implications, Conditional on Joint Completion of Studies A--G}

We separate the conditional implications from the present pilot
evidence by stating them only as conditionals. \emph{If} the program
of \S\ref{sec:requiredwork} is jointly completed and returns positive
results, then several consequences follow. The SAE feature dictionary
\cite{bricken2023monosemanticity, cunningham2024sparse,
templeton2024scaling, gao2024scaling, lieberum2024gemma} becomes
interpretable as the empirical recovery of bases for approximately
invariant subspaces of trained transformer forward operators.
Cross-model transfer of mechanistic understanding becomes geometrically
tractable via the recovered Procrustes transformation. The Platonic
Representation Hypothesis \cite{huh2024platonic} acquires a
mathematical structure that explains why convergence should occur and
what its limits are. The substrate-independence intuition descending
from Marr \cite{marr1982vision} acquires sharp empirical
operationalization on language-model substrates.

None of these consequences is established by the pilot. We list them
to motivate the studies the pilot does not perform.

\subsection{Authorship Note}

This paper is the work of one human author and one language model in
extended dialogue. The model has no continuity across sessions; the
human has the operator-theoretic intuition the paper articulates and
holds responsibility for the editorial and methodological judgments.
Neither party, alone, would have produced this paper. The question of
how authorship should be attributed in such a configuration is itself
a question the field will need to answer. The author block makes our
choice explicit; we recognize it is contestable.

\subsection*{Closing Note}

The pilot is what it is: a representational-alignment result against
the row-permutation null, with the gaps to the strong claim explicitly
catalogued and a research program specified to close them. We have
attempted to report the result honestly enough that a reader can
distinguish what the pilot supports from what the operator-theoretic
program would, if jointly successful, support. The next step is
Study~A: direct measurement of the invariance leakage statistic the
paper defines on page two. Until that statistic is computed and
reported, the operator-theoretic interpretation is a hypothesis the
pilot motivates rather than a finding the pilot confirms.
